\begin{lemma}[MellinOfSmooth1a]
\uses{Smooth1}
\label{MellinOfSmooth1a}
\lean{MellinOfSmooth1a}
% at /home/runner/work/trough/trough/library/PrimeNumberTheoremAnd/MellinCalculus.lean:1198.0-1321.8
\leanok
Fix  $\epsilon>0$. Then the Mellin transform of $\widetilde{1_{\epsilon}}$ is
  $$\mathcal{M}(\widetilde{1_{\epsilon}})(s) =
  \frac{1}{s}\left(\mathcal{M}(\nu)\left(\epsilon s\right)\right).$$
\end{lemma}
\begin{proof}
\uses{MellinOf1,MellinOfDeltaSpike,DeltaSpike,MellinConvolution,MellinConvolutionTransform,MellinConvolutionSymmetric}
\leanok
By Definition \ref{Smooth1},
  $$
    \mathcal M(\widetilde{1_\epsilon})(s)
    =\mathcal M(1_{(0,1]}\ast\nu_\epsilon)(s)
    .
  $$
  We wish to apply Theorem \ref{MellinConvolutionTransform}.
  To do so, we must prove that
  $$
    (x,y)\mapsto 1_{(0,1]}(y)\nu_\epsilon(x/y)/y
  $$
  is integrable on $[0,\infty)^2$.
  It is actually easier to do this for the convolution: $\nu_\epsilon\ast 1_{(0,1]}$, so we use
  Lemma \ref{MellinConvolutionSymmetric}: for $x\neq0$,
  $$
    1_{(0,1]}\ast\nu_\epsilon(x)=\nu_\epsilon\ast 1_{(0,1]}(x)
    .
  $$
  Now, for $x=0$, both sides of the equation are 0, so the equation also holds for $x=0$.
  Therefore,
  $$
    \mathcal M(\widetilde{1_\epsilon})(s)
    =\mathcal M(\nu_\epsilon\ast 1_{(0,1]})(s)
    .
  $$
  Now,
  $$
    (x,y)\mapsto \nu_\epsilon(y)1_{(0,1]}(x/y)\frac{x^{s-1}}y
  $$
  has compact support that is bounded away from $y=0$ (specifically
  $y\in[2^{-\epsilon},2^\epsilon]$ and $x\in(0,y]$), so it is integrable.
  We can thus apply Theorem \ref{MellinConvolutionTransform} and find
  $$
    \mathcal M(\widetilde{1_\epsilon})(s)
    =\mathcal M(\nu_\epsilon)(s)\mathcal M(1_{(0,1]})(s)
    .
  $$
  By Lemmas \ref{MellinOf1} and \ref{MellinOfDeltaSpike},
  $$
    \mathcal M(\widetilde{1_\epsilon})(s)
    =\frac1s\mathcal M(\nu)(\epsilon s)
    .
  $$
\end{proof}
